% Copyright 2009 by Tomasz Mazur
%
% This file may be distributed and/or modified in all ways.


\documentclass[xcolor=pdftex,t,11pt]{beamer}

%%%%%%%%%%%%%%%%%%%%%%%%%%%%%%%%%%
%       SET OPTIONS BELOW        %
%%%%%%%%%%%%%%%%%%%%%%%%%%%%%%%%%%

\usetheme[
% Toggle showing page counter
pagecounter=true,
%
% String to be used between the current page and the
% total page count, e.g. of, /, from, etc.
pageofpages=of,
%
% Defines the shape of bullet points. Available options: circle, square
bullet=circle,
%
% Show a line below the frame title. 
titleline=true,
%
% Set the style of the title page (true for fancy, false for standard)
alternativetitlepage=true,
%
% Institution logo for fancy title page.
% Comment out to remove the logo from the title page.
% IMPORTANT: THERE IS A BUG IN SOME VERSIONS OF PDFLATEX AND FONTS
% ON THE LOGOS ARE NOT RENDERED PROPERLY. IN SUCH A CASE ADD `2` 
% TO THE NAME OF THE LOGO, E.G. comlab2 INSTEAD OF comlab
%titlepagelogo=images/titlepage/ou, % Phiri --October 30, 2012. Commented out
%
% Department footer logo for fancy title page
% Comment out to remove the logo from the footer of the title page/
% IMPORTANT: THERE IS A BUG IN SOME VERSIONS OF PDFLATEX AND FONTS
% ON THE LOGOS ARE NOT RENDERED PROPERLY. IN SUCH A CASE ADD `2` 
% TO THE NAME OF THE LOGO, E.G. comlab2 INSTEAD OF comlab
%titlepagefooterlogo=images/titlepage/comlab,
%
% Institution/department logo for ordinary slides
% Comment this line out to remove the logo from all the pages.
% Available logos are: ou, comlab, comlabinline, comlabou
% IMPORTANT: THERE IS A BUG IN SOME VERSIONS OF PDFLATEX AND FONTS
% ON THE LOGOS ARE NOT RENDERED PROPERLY. IN SUCH A CASE ADD `2` 
% TO THE NAME OF THE LOGO, E.G. comlab2 INSTEAD OF comlab
%ordinarypageslogo=images/comlab, % phiri --October 30, 2012. Commented out
%
%
% Add watermark in the bottom right corner
%watermark=<filename>,
%
% Set the height of the watermark.
%watermarkheight=100pt,
%
% The watermark image is 4 times bigger than watermarkheight.
%watermarkheightmult=4,
]{Torino}

% Select color theme. Available options are:
% mininmal, greenandblue, blue, red
\usecolortheme{minimal}

%Select different font themes.Available options are:
% default, serif, structurebold, structureitalicserif, structuresmallcapsserif
\usefonttheme{structurebold}


%%%%%%%%%%%%%%%%%%%%%%%%%%%%%%%%%%
%       PRESENTATION INFO        %
%%%%%%%%%%%%%%%%%%%%%%%%%%%%%%%%%%

\author{Lighton Phiri \and Kyle Williams \and Miles Robinson\\Stuart Hammar \and Hussein Suleman}
\title{Bonolo\symbolfootnote[2, frame]{\tiny Sotho word meaning easy.}}
\subtitle{A General Digital Library System for File-based Collections}
\institute{Digital Libraries Laboratory\\Department of Computer Science\\University of Cape Town}
\date{November 13, 2012}

% Footer symbols
\long\def\symbolfootnote[#1]#2{\begingroup%
\def\thefootnote{\fnsymbol{footnote}}\footnote[#1]{#2}\endgroup}

%
%Replace default bullets to square bullets
%
\setbeamertemplate{items}[square]

% Longtable lets you have tables that span multiple pages.
\usepackage{longtable}

%%% For tables
\usepackage{multirow}

% Rotating table headers/labels
\usepackage{rotating}

%package to enable tikZ
\usepackage{tikz}

% Footer collision with page numbers
\addtobeamertemplate{footnote}{}{\vspace{2ex}}

\usepackage{Sweave}
\begin{document}



%%%%%%%%%%%%%%%%%%%%%%%%%%%%%%%%%%
%       SLIDE DEFINITIONS        %
%%%%%%%%%%%%%%%%%%%%%%%%%%%%%%%%%%

\begin{frame}[plain]
	\titlepage
\end{frame}

\begin{frame}{Outline}
	\tableofcontents
\end{frame}


\section{Introduction}


\subsection{What is this?}

\begin{frame}[fragile]
\frametitle{Introduction}
\begin{itemize}
\item \alert{Creating presentations using \LaTeX{} is straightforward...}
\item ...with Beamer, a class for creating slides
\item should be already installed with most \LaTeX{} distributions, but can be obtained from \verb!http://latex-beamer.sourceforge.net/!
\item Beamer documentation available from \verb!	http://www.ctan.org/tex-archive/macros/latex/! \verb!	contrib/beamer/doc/beameruserguide.pdf!
\item This is a modification of Marco Barision's Torino theme 
\item It aims to produce slides that are \emph{pretty}, but easily \emph{readable} and with \emph{large content area}
\item Most of standard Beamer commands are supported
\end{itemize}
\end{frame}

\subsection{Creating presentation overview}

\begin{frame}[fragile]
\frametitle{Status Quo - DL Architectural Frameworks}
\begin{itemize}
\item Warwick Framework
\item Khan/Wilensky Framework
\item DELOS DL Reference Model
\item 5S Framework
\end{itemize}
\end{frame}


\section{Modifying themes, colours and fonts}


\begin{frame}[fragile]
\frametitle{Status Quo - DL Software}
\begin{itemize}
\item Splash screen with prominent DL FLOSS software
\begin{itemize}
	\item Common elements
	\item Fundemental principles
	\item Suitability
	\item Cost implications
\end{itemize}
\item Additional information $\cdots$
\item Additional information $\cdots$
\end{itemize}
\end{frame}

\begin{frame}[fragile]
\frametitle{Motivation}
\begin{itemize}
\item There are five font themes: 
\begin{itemize}
	\item \verb!default! (sans serif)
	\item \verb!serif! (used for this presentation) 
	\item \verb!structurebold! (titles, headlines, etc.\ are typeset in a bold font)
	\item \verb!structureitalicserif! (titles, headlines, etc.\ are typeset in an italics serif font)
	\item \verb!structuresmallcapsserif! (titles, headlines, etc.\ are typeset in a small caps serif font)
\end{itemize}
\item Change the document-wise font size to 8, 9 , 10, 11 (default), 12, 14, 17 or 20 points in the options of \verb!\documentclass!, e.g. \verb!\documentclass[12pt]{beamer}!
\item Colour text using \verb!\textcolor{<<colour>>}{<<text>>}!
\item The \verb!\alert{<<text>>}! command colours text red
\end{itemize}
\end{frame}

\subsection{Using boxes and images}

\begin{frame}[fragile]
\frametitle{Design Principles}
\begin{itemize}
\item Exploratory Study
\begin{itemize}
\item Objective of study
\item 
\end{itemize}
\item Grounded Theory approach
\begin{itemize}
\item Data collection involved meta-analysis of twelve (12) software tools
\item Data analysis yielded eight (8) design principles
\end{itemize}
\item Outcome of study phase yieded eight (8) design principles
\end{itemize}
\end{frame}

\begin{frame}[fragile]
\frametitle{Design Principles}
\begin{itemize}
\item Hardware and/or software platform independence
\item Heterogeneous object, metadata and service integration
\item Support for community and international standards
\item Flexible design to facilitate extensibility
\item Minimalist design approach
\item Simplified preservation process
\item Structured organisation of data
\item Design for least possible resources
\end{itemize}

\end{frame}


\begin{frame}[fragile]
\frametitle{Prototype Implementation - Bonolo}
\begin{itemize}
\item Architectural diagram goes here
\end{itemize}
\begin{figure}
%\includegraphics[height=100pt]{images/titlepage/ou}
%\caption{Oxford University logo}
\end{figure}
\end{frame}

\begin{frame}[fragile]
\frametitle{End User Interface - Item View}
\begin{figure}
\includegraphics{images/icadl2012_bonolo_end_user_interface.png}
%\caption{End User Interface}
\end{figure}
\end{frame}

\begin{frame}[fragile]
\frametitle{Curator Interface - Item View}
\begin{figure}
\includegraphics{images/bonolo-ci-item.png}
%\caption{Curator Interface}
\end{figure}
\end{frame}

\begin{frame}[fragile]
\frametitle{System Design - Repository}
\begin{itemize}
\item Architectural diagram goes here
\end{itemize}
\begin{figure}
%\includegraphics[height=100pt]{images/titlepage/ou}
%\caption{Oxford University logo}
\end{figure}
\end{frame}


\section{Overlays}


\begin{frame}[fragile]
\frametitle{Simple overlays}
\begin{itemize}
\item<1-> Use \verb!\begin{itemize} \item<x-> \end{itemize}! to display bullet points incrementally
\item<2-> Alternatively, use the \verb!\pause! command, which displays contents of the slide up to the first marker, then up to the second marker, etc.
\end{itemize}
\pause
\begin{table}
\begin{tabular}{lcccc}
        & A & B & C & D \\\hline
  X     & 1 & 2 & 3 & 4 \pause\\
  Y     & 3 & 4 & 5 & 6 \pause\\
  Z     & 5 & 6 & 7 & 8
\end{tabular}
\end{table}
\end{frame}

\begin{frame}[fragile]
\frametitle{End User Interface UX - Lab Experiment}
\begin{itemize}
\item Objective
\item Target Group
\item Approach
\end{itemize}
\end{frame}


\begin{frame}[fragile]
\frametitle{End User Interface UX - Experiment Results}
\centering
% Created by tikzDevice version 0.6.2-92-0ad2792 on 2012-11-10 03:17:29
% !TEX encoding = UTF-8 Unicode
\begin{tikzpicture}[x=1pt,y=1pt]
\definecolor[named]{fillColor}{rgb}{1.00,1.00,1.00}
\path[use as bounding box,fill=fillColor,fill opacity=0.00] (0,0) rectangle (332.44,195.13);
\begin{scope}
\path[clip] (  0.00,  0.00) rectangle (332.44,195.13);
\definecolor[named]{drawColor}{rgb}{0.00,0.00,0.00}
\definecolor[named]{fillColor}{rgb}{1.00,0.00,0.00}

\path[draw=drawColor,line width= 0.4pt,line join=round,line cap=round,fill=fillColor] ( 38.71, 28.80) rectangle ( 45.15,156.10);
\definecolor[named]{fillColor}{rgb}{1.00,0.86,0.00}

\path[draw=drawColor,line width= 0.4pt,line join=round,line cap=round,fill=fillColor] ( 45.15, 28.80) rectangle ( 51.59,175.20);
\definecolor[named]{fillColor}{rgb}{0.29,1.00,0.00}

\path[draw=drawColor,line width= 0.4pt,line join=round,line cap=round,fill=fillColor] ( 51.59, 28.80) rectangle ( 58.02,172.02);
\definecolor[named]{fillColor}{rgb}{0.00,1.00,0.57}

\path[draw=drawColor,line width= 0.4pt,line join=round,line cap=round,fill=fillColor] ( 58.02, 28.80) rectangle ( 64.46,178.38);
\definecolor[named]{fillColor}{rgb}{0.00,0.57,1.00}

\path[draw=drawColor,line width= 0.4pt,line join=round,line cap=round,fill=fillColor] ( 64.46, 28.80) rectangle ( 70.90,184.75);
\definecolor[named]{fillColor}{rgb}{0.29,0.00,1.00}

\path[draw=drawColor,line width= 0.4pt,line join=round,line cap=round,fill=fillColor] ( 70.90, 28.80) rectangle ( 77.33,178.38);
\definecolor[named]{fillColor}{rgb}{1.00,0.00,0.86}

\path[draw=drawColor,line width= 0.4pt,line join=round,line cap=round,fill=fillColor] ( 77.33, 28.80) rectangle ( 83.77,181.56);
\definecolor[named]{fillColor}{rgb}{1.00,0.00,0.00}

\path[draw=drawColor,line width= 0.4pt,line join=round,line cap=round,fill=fillColor] (106.30, 28.80) rectangle (112.74, 28.80);
\definecolor[named]{fillColor}{rgb}{1.00,0.86,0.00}

\path[draw=drawColor,line width= 0.4pt,line join=round,line cap=round,fill=fillColor] (112.74, 28.80) rectangle (119.17,168.83);
\definecolor[named]{fillColor}{rgb}{0.29,1.00,0.00}

\path[draw=drawColor,line width= 0.4pt,line join=round,line cap=round,fill=fillColor] (119.17, 28.80) rectangle (125.61,168.83);
\definecolor[named]{fillColor}{rgb}{0.00,1.00,0.57}

\path[draw=drawColor,line width= 0.4pt,line join=round,line cap=round,fill=fillColor] (125.61, 28.80) rectangle (132.05,172.02);
\definecolor[named]{fillColor}{rgb}{0.00,0.57,1.00}

\path[draw=drawColor,line width= 0.4pt,line join=round,line cap=round,fill=fillColor] (132.05, 28.80) rectangle (138.48, 28.80);
\definecolor[named]{fillColor}{rgb}{0.29,0.00,1.00}

\path[draw=drawColor,line width= 0.4pt,line join=round,line cap=round,fill=fillColor] (138.48, 28.80) rectangle (144.92, 28.80);
\definecolor[named]{fillColor}{rgb}{1.00,0.00,0.86}

\path[draw=drawColor,line width= 0.4pt,line join=round,line cap=round,fill=fillColor] (144.92, 28.80) rectangle (151.36, 28.80);
\definecolor[named]{fillColor}{rgb}{1.00,0.00,0.00}

\path[draw=drawColor,line width= 0.4pt,line join=round,line cap=round,fill=fillColor] (173.89, 28.80) rectangle (180.32,168.83);
\definecolor[named]{fillColor}{rgb}{1.00,0.86,0.00}

\path[draw=drawColor,line width= 0.4pt,line join=round,line cap=round,fill=fillColor] (180.32, 28.80) rectangle (186.76,168.83);
\definecolor[named]{fillColor}{rgb}{0.29,1.00,0.00}

\path[draw=drawColor,line width= 0.4pt,line join=round,line cap=round,fill=fillColor] (186.76, 28.80) rectangle (193.20, 28.80);
\definecolor[named]{fillColor}{rgb}{0.00,1.00,0.57}

\path[draw=drawColor,line width= 0.4pt,line join=round,line cap=round,fill=fillColor] (193.20, 28.80) rectangle (199.63,172.02);
\definecolor[named]{fillColor}{rgb}{0.00,0.57,1.00}

\path[draw=drawColor,line width= 0.4pt,line join=round,line cap=round,fill=fillColor] (199.63, 28.80) rectangle (206.07,181.56);
\definecolor[named]{fillColor}{rgb}{0.29,0.00,1.00}

\path[draw=drawColor,line width= 0.4pt,line join=round,line cap=round,fill=fillColor] (206.07, 28.80) rectangle (212.51, 28.80);
\definecolor[named]{fillColor}{rgb}{1.00,0.00,0.86}

\path[draw=drawColor,line width= 0.4pt,line join=round,line cap=round,fill=fillColor] (212.51, 28.80) rectangle (218.94, 28.80);
\definecolor[named]{fillColor}{rgb}{1.00,0.00,0.00}

\path[draw=drawColor,line width= 0.4pt,line join=round,line cap=round,fill=fillColor] (241.47, 28.80) rectangle (247.91,172.02);
\definecolor[named]{fillColor}{rgb}{1.00,0.86,0.00}

\path[draw=drawColor,line width= 0.4pt,line join=round,line cap=round,fill=fillColor] (247.91, 28.80) rectangle (254.35,168.83);
\definecolor[named]{fillColor}{rgb}{0.29,1.00,0.00}

\path[draw=drawColor,line width= 0.4pt,line join=round,line cap=round,fill=fillColor] (254.35, 28.80) rectangle (260.78,172.02);
\definecolor[named]{fillColor}{rgb}{0.00,1.00,0.57}

\path[draw=drawColor,line width= 0.4pt,line join=round,line cap=round,fill=fillColor] (260.78, 28.80) rectangle (267.22, 28.80);
\definecolor[named]{fillColor}{rgb}{0.00,0.57,1.00}

\path[draw=drawColor,line width= 0.4pt,line join=round,line cap=round,fill=fillColor] (267.22, 28.80) rectangle (273.66, 28.80);
\definecolor[named]{fillColor}{rgb}{0.29,0.00,1.00}

\path[draw=drawColor,line width= 0.4pt,line join=round,line cap=round,fill=fillColor] (273.66, 28.80) rectangle (280.09, 28.80);
\definecolor[named]{fillColor}{rgb}{1.00,0.00,0.86}

\path[draw=drawColor,line width= 0.4pt,line join=round,line cap=round,fill=fillColor] (280.09, 28.80) rectangle (286.53, 28.80);
\end{scope}
\begin{scope}
\path[clip] (  0.00,  0.00) rectangle (332.44,195.13);
\definecolor[named]{drawColor}{rgb}{0.00,0.00,0.00}

\node[text=drawColor,anchor=base,inner sep=0pt, outer sep=0pt, scale=  0.60] at ( 61.24, 15.84) {Search\& Browse};

\node[text=drawColor,anchor=base,inner sep=0pt, outer sep=0pt, scale=  0.60] at (128.83, 15.84) {Navigability};

\node[text=drawColor,anchor=base,inner sep=0pt, outer sep=0pt, scale=  0.60] at (196.41, 15.84) {Similarity};

\node[text=drawColor,anchor=base,inner sep=0pt, outer sep=0pt, scale=  0.60] at (264.00, 15.84) {General};
\end{scope}
\begin{scope}
\path[clip] (  0.00,  0.00) rectangle (332.44,195.13);
\definecolor[named]{drawColor}{rgb}{0.00,0.00,0.00}

\node[text=drawColor,anchor=base,inner sep=0pt, outer sep=0pt, scale=  0.60] at (162.62,  1.44) {Average UX Ratings};

\node[text=drawColor,rotate= 90.00,anchor=base,inner sep=0pt, outer sep=0pt, scale=  0.60] at (  5.76,108.36) {Survey Participants};
\end{scope}
\begin{scope}
\path[clip] (  0.00,  0.00) rectangle (332.44,195.13);
\definecolor[named]{drawColor}{rgb}{0.00,0.00,0.00}

\path[draw=drawColor,line width= 0.4pt,line join=round,line cap=round] ( 28.80, 28.80) -- ( 28.80,187.93);

\path[draw=drawColor,line width= 0.4pt,line join=round,line cap=round] ( 28.80, 28.80) -- ( 25.20, 28.80);

\path[draw=drawColor,line width= 0.4pt,line join=round,line cap=round] ( 28.80, 60.63) -- ( 25.20, 60.63);

\path[draw=drawColor,line width= 0.4pt,line join=round,line cap=round] ( 28.80, 92.45) -- ( 25.20, 92.45);

\path[draw=drawColor,line width= 0.4pt,line join=round,line cap=round] ( 28.80,124.28) -- ( 25.20,124.28);

\path[draw=drawColor,line width= 0.4pt,line join=round,line cap=round] ( 28.80,156.10) -- ( 25.20,156.10);

\path[draw=drawColor,line width= 0.4pt,line join=round,line cap=round] ( 28.80,187.93) -- ( 25.20,187.93);

\node[text=drawColor,anchor=base east,inner sep=0pt, outer sep=0pt, scale=  0.60] at ( 21.60, 26.73) {0};

\node[text=drawColor,anchor=base east,inner sep=0pt, outer sep=0pt, scale=  0.60] at ( 21.60, 58.56) {1};

\node[text=drawColor,anchor=base east,inner sep=0pt, outer sep=0pt, scale=  0.60] at ( 21.60, 90.39) {2};

\node[text=drawColor,anchor=base east,inner sep=0pt, outer sep=0pt, scale=  0.60] at ( 21.60,122.21) {3};

\node[text=drawColor,anchor=base east,inner sep=0pt, outer sep=0pt, scale=  0.60] at ( 21.60,154.04) {4};

\node[text=drawColor,anchor=base east,inner sep=0pt, outer sep=0pt, scale=  0.60] at ( 21.60,185.86) {5};
\end{scope}
\begin{scope}
\path[clip] (  0.00,  0.00) rectangle (332.44,195.13);
\definecolor[named]{drawColor}{rgb}{0.00,0.00,0.00}

\path[draw=drawColor,line width= 0.4pt,line join=round,line cap=round] (280.09,187.93) rectangle (323.16,130.33);
\definecolor[named]{fillColor}{rgb}{1.00,0.00,0.00}

\path[draw=drawColor,line width= 0.4pt,line join=round,line cap=round,fill=fillColor] (285.49,182.53) rectangle (289.81,178.93);
\definecolor[named]{fillColor}{rgb}{1.00,0.86,0.00}

\path[draw=drawColor,line width= 0.4pt,line join=round,line cap=round,fill=fillColor] (285.49,175.33) rectangle (289.81,171.73);
\definecolor[named]{fillColor}{rgb}{0.29,1.00,0.00}

\path[draw=drawColor,line width= 0.4pt,line join=round,line cap=round,fill=fillColor] (285.49,168.13) rectangle (289.81,164.53);
\definecolor[named]{fillColor}{rgb}{0.00,1.00,0.57}

\path[draw=drawColor,line width= 0.4pt,line join=round,line cap=round,fill=fillColor] (285.49,160.93) rectangle (289.81,157.33);
\definecolor[named]{fillColor}{rgb}{0.00,0.57,1.00}

\path[draw=drawColor,line width= 0.4pt,line join=round,line cap=round,fill=fillColor] (285.49,153.73) rectangle (289.81,150.13);
\definecolor[named]{fillColor}{rgb}{0.29,0.00,1.00}

\path[draw=drawColor,line width= 0.4pt,line join=round,line cap=round,fill=fillColor] (285.49,146.53) rectangle (289.81,142.93);
\definecolor[named]{fillColor}{rgb}{1.00,0.00,0.86}

\path[draw=drawColor,line width= 0.4pt,line join=round,line cap=round,fill=fillColor] (285.49,139.33) rectangle (289.81,135.73);

\node[text=drawColor,anchor=base west,inner sep=0pt, outer sep=0pt, scale=  0.60] at (295.21,178.66) {Intuitive};

\node[text=drawColor,anchor=base west,inner sep=0pt, outer sep=0pt, scale=  0.60] at (295.21,171.46) {Simple};

\node[text=drawColor,anchor=base west,inner sep=0pt, outer sep=0pt, scale=  0.60] at (295.21,164.26) {Satisfied};

\node[text=drawColor,anchor=base west,inner sep=0pt, outer sep=0pt, scale=  0.60] at (295.21,157.06) {Expected};

\node[text=drawColor,anchor=base west,inner sep=0pt, outer sep=0pt, scale=  0.60] at (295.21,149.86) {Response};

\node[text=drawColor,anchor=base west,inner sep=0pt, outer sep=0pt, scale=  0.60] at (295.21,142.66) {Effective};

\node[text=drawColor,anchor=base west,inner sep=0pt, outer sep=0pt, scale=  0.60] at (295.21,135.46) {Beneficial};
\end{scope}
\end{tikzpicture}

\end{frame}

\begin{frame}[fragile]
\frametitle{Curator Interface UX - Online Experiment}
\begin{itemize}
\item Rationale
\begin{itemize}
\item User experience working with curation features
\end{itemize}
\item Target Group
\begin{itemize}
\item Individuals with no experience working with DL tools
\item Social networking site recruitment\footnote{\tiny Facebook, Twitter,
Google$+$}
\end{itemize}
\item Approach
\begin{itemize}
\item Intrinsic Motivation Inventory
\item Five (5) minute 'HOWTO' screencast
\item Curation tasks with two datasets
\item Online questionnaire
\end{itemize}
\end{itemize}
\end{frame}


\begin{frame}[fragile]
\frametitle{Curator Interface UX - Experiment Results}
\centering
% Created by tikzDevice version 0.6.2-92-0ad2792 on 2012-11-10 04:50:24
% !TEX encoding = UTF-8 Unicode
\begin{tikzpicture}[x=1pt,y=1pt]
\definecolor[named]{fillColor}{rgb}{1.00,1.00,1.00}
\path[use as bounding box,fill=fillColor,fill opacity=0.00] (0,0) rectangle (332.44,195.13);
\begin{scope}
\path[clip] (  0.00,  0.00) rectangle (332.44,195.13);
\definecolor[named]{drawColor}{rgb}{0.00,0.00,0.00}
\definecolor[named]{fillColor}{rgb}{1.00,0.00,0.00}

\path[draw=drawColor,line width= 0.4pt,line join=round,line cap=round,fill=fillColor] ( 36.85, 28.80) rectangle ( 96.01, 41.53);
\definecolor[named]{fillColor}{rgb}{0.80,1.00,0.00}

\path[draw=drawColor,line width= 0.4pt,line join=round,line cap=round,fill=fillColor] ( 36.85, 41.53) rectangle ( 96.01, 66.99);
\definecolor[named]{fillColor}{rgb}{0.00,1.00,0.40}

\path[draw=drawColor,line width= 0.4pt,line join=round,line cap=round,fill=fillColor] ( 36.85, 66.99) rectangle ( 96.01,137.01);
\definecolor[named]{fillColor}{rgb}{0.00,0.40,1.00}

\path[draw=drawColor,line width= 0.4pt,line join=round,line cap=round,fill=fillColor] ( 36.85,137.01) rectangle ( 96.01,162.47);
\definecolor[named]{fillColor}{rgb}{0.80,0.00,1.00}

\path[draw=drawColor,line width= 0.4pt,line join=round,line cap=round,fill=fillColor] ( 36.85,162.47) rectangle ( 96.01,168.83);
\definecolor[named]{fillColor}{rgb}{1.00,0.00,0.00}

\path[draw=drawColor,line width= 0.4pt,line join=round,line cap=round,fill=fillColor] (107.84, 28.80) rectangle (167.00, 47.90);
\definecolor[named]{fillColor}{rgb}{0.80,1.00,0.00}

\path[draw=drawColor,line width= 0.4pt,line join=round,line cap=round,fill=fillColor] (107.84, 47.90) rectangle (167.00,143.37);
\definecolor[named]{fillColor}{rgb}{0.00,1.00,0.40}

\path[draw=drawColor,line width= 0.4pt,line join=round,line cap=round,fill=fillColor] (107.84,143.37) rectangle (167.00,156.10);
\definecolor[named]{fillColor}{rgb}{0.00,0.40,1.00}

\path[draw=drawColor,line width= 0.4pt,line join=round,line cap=round,fill=fillColor] (107.84,156.10) rectangle (167.00,162.47);
\definecolor[named]{fillColor}{rgb}{0.80,0.00,1.00}

\path[draw=drawColor,line width= 0.4pt,line join=round,line cap=round,fill=fillColor] (107.84,162.47) rectangle (167.00,168.83);
\definecolor[named]{fillColor}{rgb}{1.00,0.00,0.00}

\path[draw=drawColor,line width= 0.4pt,line join=round,line cap=round,fill=fillColor] (178.83, 28.80) rectangle (238.00, 79.72);
\definecolor[named]{fillColor}{rgb}{0.80,1.00,0.00}

\path[draw=drawColor,line width= 0.4pt,line join=round,line cap=round,fill=fillColor] (178.83, 79.72) rectangle (238.00,149.74);
\definecolor[named]{fillColor}{rgb}{0.00,1.00,0.40}

\path[draw=drawColor,line width= 0.4pt,line join=round,line cap=round,fill=fillColor] (178.83,149.74) rectangle (238.00,156.10);
\definecolor[named]{fillColor}{rgb}{0.00,0.40,1.00}

\path[draw=drawColor,line width= 0.4pt,line join=round,line cap=round,fill=fillColor] (178.83,156.10) rectangle (238.00,168.83);
\definecolor[named]{fillColor}{rgb}{0.80,0.00,1.00}

\path[draw=drawColor,line width= 0.4pt,line join=round,line cap=round,fill=fillColor] (178.83,168.83) rectangle (238.00,168.83);
\end{scope}
\begin{scope}
\path[clip] (  0.00,  0.00) rectangle (332.44,195.13);
\definecolor[named]{drawColor}{rgb}{0.00,0.00,0.00}

\node[text=drawColor,anchor=base,inner sep=0pt, outer sep=0pt, scale=  0.60] at ( 66.43, 15.84) {Interest/Enjoyment};

\node[text=drawColor,anchor=base,inner sep=0pt, outer sep=0pt, scale=  0.60] at (137.42, 15.84) {Perceived Competence};

\node[text=drawColor,anchor=base,inner sep=0pt, outer sep=0pt, scale=  0.60] at (208.42, 15.84) {Value/Usefulness};
\end{scope}
\begin{scope}
\path[clip] (  0.00,  0.00) rectangle (332.44,195.13);
\definecolor[named]{drawColor}{rgb}{0.00,0.00,0.00}

\node[text=drawColor,anchor=base,inner sep=0pt, outer sep=0pt, scale=  0.60] at (137.42,  1.44) {Average UX Ratings};

\node[text=drawColor,rotate= 90.00,anchor=base,inner sep=0pt, outer sep=0pt, scale=  0.60] at (  5.76,108.36) {Survey Participants};
\end{scope}
\begin{scope}
\path[clip] (  0.00,  0.00) rectangle (332.44,195.13);
\definecolor[named]{drawColor}{rgb}{0.00,0.00,0.00}

\path[draw=drawColor,line width= 0.4pt,line join=round,line cap=round] ( 28.80, 28.80) -- ( 28.80,187.93);

\path[draw=drawColor,line width= 0.4pt,line join=round,line cap=round] ( 28.80, 28.80) -- ( 25.20, 28.80);

\path[draw=drawColor,line width= 0.4pt,line join=round,line cap=round] ( 28.80, 60.63) -- ( 25.20, 60.63);

\path[draw=drawColor,line width= 0.4pt,line join=round,line cap=round] ( 28.80, 92.45) -- ( 25.20, 92.45);

\path[draw=drawColor,line width= 0.4pt,line join=round,line cap=round] ( 28.80,124.28) -- ( 25.20,124.28);

\path[draw=drawColor,line width= 0.4pt,line join=round,line cap=round] ( 28.80,156.10) -- ( 25.20,156.10);

\path[draw=drawColor,line width= 0.4pt,line join=round,line cap=round] ( 28.80,187.93) -- ( 25.20,187.93);

\node[text=drawColor,anchor=base east,inner sep=0pt, outer sep=0pt, scale=  0.60] at ( 21.60, 26.73) {0};

\node[text=drawColor,anchor=base east,inner sep=0pt, outer sep=0pt, scale=  0.60] at ( 21.60, 58.56) {5};

\node[text=drawColor,anchor=base east,inner sep=0pt, outer sep=0pt, scale=  0.60] at ( 21.60, 90.39) {10};

\node[text=drawColor,anchor=base east,inner sep=0pt, outer sep=0pt, scale=  0.60] at ( 21.60,122.21) {15};

\node[text=drawColor,anchor=base east,inner sep=0pt, outer sep=0pt, scale=  0.60] at ( 21.60,154.04) {20};

\node[text=drawColor,anchor=base east,inner sep=0pt, outer sep=0pt, scale=  0.60] at ( 21.60,185.86) {25};
\end{scope}
\begin{scope}
\path[clip] (  0.00,  0.00) rectangle (332.44,195.13);
\definecolor[named]{drawColor}{rgb}{0.00,0.00,0.00}

\path[draw=drawColor,line width= 0.4pt,line join=round,line cap=round] (258.70,181.56) rectangle (323.00,138.36);
\definecolor[named]{fillColor}{rgb}{1.00,0.00,0.00}

\path[draw=drawColor,line width= 0.4pt,line join=round,line cap=round,fill=fillColor] (264.10,176.16) rectangle (268.42,172.56);
\definecolor[named]{fillColor}{rgb}{0.80,1.00,0.00}

\path[draw=drawColor,line width= 0.4pt,line join=round,line cap=round,fill=fillColor] (264.10,168.96) rectangle (268.42,165.36);
\definecolor[named]{fillColor}{rgb}{0.00,1.00,0.40}

\path[draw=drawColor,line width= 0.4pt,line join=round,line cap=round,fill=fillColor] (264.10,161.76) rectangle (268.42,158.16);
\definecolor[named]{fillColor}{rgb}{0.00,0.40,1.00}

\path[draw=drawColor,line width= 0.4pt,line join=round,line cap=round,fill=fillColor] (264.10,154.56) rectangle (268.42,150.96);
\definecolor[named]{fillColor}{rgb}{0.80,0.00,1.00}

\path[draw=drawColor,line width= 0.4pt,line join=round,line cap=round,fill=fillColor] (264.10,147.36) rectangle (268.42,143.76);

\node[text=drawColor,anchor=base west,inner sep=0pt, outer sep=0pt, scale=  0.60] at (273.82,172.30) {Strongly Agree};

\node[text=drawColor,anchor=base west,inner sep=0pt, outer sep=0pt, scale=  0.60] at (273.82,165.10) {Agree};

\node[text=drawColor,anchor=base west,inner sep=0pt, outer sep=0pt, scale=  0.60] at (273.82,157.90) {Neutral};

\node[text=drawColor,anchor=base west,inner sep=0pt, outer sep=0pt, scale=  0.60] at (273.82,150.70) {Disagree};

\node[text=drawColor,anchor=base west,inner sep=0pt, outer sep=0pt, scale=  0.60] at (273.82,143.50) {Strongly Disagree};
\end{scope}
\end{tikzpicture}

\end{frame}


\section{What else is possible}


\subsection{Making most of the Beamer class}

\begin{frame}[fragile]
\frametitle{Repository Performance - Lab Experiment}
\begin{itemize}
\item Objective
\item Target Group
\item Approach
\end{itemize}
\end{frame}


\begin{frame}[fragile]
\frametitle{Repository Performance - Experiment Results}
\centering
% Created by tikzDevice version 0.6.2-92-0ad2792 on 2012-11-10 03:17:14
% !TEX encoding = UTF-8 Unicode
\begin{tikzpicture}[x=1pt,y=1pt]
\definecolor[named]{fillColor}{rgb}{1.00,1.00,1.00}
\path[use as bounding box,fill=fillColor,fill opacity=0.00] (0,0) rectangle (332.44,195.13);
\begin{scope}
\path[clip] (  0.00,  0.00) rectangle (332.44,195.13);
\definecolor[named]{drawColor}{rgb}{0.00,0.00,1.00}

\path[draw=drawColor,line width= 0.4pt,line join=round,line cap=round] ( 38.71, 37.11) --
	(100.67, 36.87) --
	(162.62, 36.99) --
	(224.58, 37.20) --
	(286.53, 37.17);

\path[draw=drawColor,line width= 0.4pt,line join=round,line cap=round] ( 38.71, 37.11) circle (  1.35);

\path[draw=drawColor,line width= 0.4pt,line join=round,line cap=round] (100.67, 36.87) circle (  1.35);

\path[draw=drawColor,line width= 0.4pt,line join=round,line cap=round] (162.62, 36.99) circle (  1.35);

\path[draw=drawColor,line width= 0.4pt,line join=round,line cap=round] (224.58, 37.20) circle (  1.35);

\path[draw=drawColor,line width= 0.4pt,line join=round,line cap=round] (286.53, 37.17) circle (  1.35);
\end{scope}
\begin{scope}
\path[clip] (  0.00,  0.00) rectangle (332.44,195.13);
\definecolor[named]{drawColor}{rgb}{0.00,0.00,0.00}

\path[draw=drawColor,line width= 0.4pt,line join=round,line cap=round] ( 38.71, 28.80) -- (286.53, 28.80);

\path[draw=drawColor,line width= 0.4pt,line join=round,line cap=round] ( 38.71, 28.80) -- ( 38.71, 25.20);

\path[draw=drawColor,line width= 0.4pt,line join=round,line cap=round] (100.67, 28.80) -- (100.67, 25.20);

\path[draw=drawColor,line width= 0.4pt,line join=round,line cap=round] (162.62, 28.80) -- (162.62, 25.20);

\path[draw=drawColor,line width= 0.4pt,line join=round,line cap=round] (224.58, 28.80) -- (224.58, 25.20);

\path[draw=drawColor,line width= 0.4pt,line join=round,line cap=round] (286.53, 28.80) -- (286.53, 25.20);

\node[text=drawColor,anchor=base,inner sep=0pt, outer sep=0pt, scale=  0.60] at ( 38.71, 15.84) {1,024};

\node[text=drawColor,anchor=base,inner sep=0pt, outer sep=0pt, scale=  0.60] at (100.67, 15.84) {2,048};

\node[text=drawColor,anchor=base,inner sep=0pt, outer sep=0pt, scale=  0.60] at (162.62, 15.84) {4,096};

\node[text=drawColor,anchor=base,inner sep=0pt, outer sep=0pt, scale=  0.60] at (224.58, 15.84) {8,192};

\node[text=drawColor,anchor=base,inner sep=0pt, outer sep=0pt, scale=  0.60] at (286.53, 15.84) {16,384};

\path[draw=drawColor,line width= 0.4pt,line join=round,line cap=round] ( 28.80, 34.69) -- ( 28.80,182.04);

\path[draw=drawColor,line width= 0.4pt,line join=round,line cap=round] ( 28.80, 34.69) -- ( 25.20, 34.69);

\path[draw=drawColor,line width= 0.4pt,line join=round,line cap=round] ( 28.80, 64.16) -- ( 25.20, 64.16);

\path[draw=drawColor,line width= 0.4pt,line join=round,line cap=round] ( 28.80, 93.63) -- ( 25.20, 93.63);

\path[draw=drawColor,line width= 0.4pt,line join=round,line cap=round] ( 28.80,123.10) -- ( 25.20,123.10);

\path[draw=drawColor,line width= 0.4pt,line join=round,line cap=round] ( 28.80,152.57) -- ( 25.20,152.57);

\path[draw=drawColor,line width= 0.4pt,line join=round,line cap=round] ( 28.80,182.04) -- ( 25.20,182.04);

\node[text=drawColor,anchor=base east,inner sep=0pt, outer sep=0pt, scale=  0.60] at ( 21.60, 32.63) {0};

\node[text=drawColor,anchor=base east,inner sep=0pt, outer sep=0pt, scale=  0.60] at ( 21.60, 62.10) {1000};

\node[text=drawColor,anchor=base east,inner sep=0pt, outer sep=0pt, scale=  0.60] at ( 21.60, 91.56) {2000};

\node[text=drawColor,anchor=base east,inner sep=0pt, outer sep=0pt, scale=  0.60] at ( 21.60,121.03) {3000};

\node[text=drawColor,anchor=base east,inner sep=0pt, outer sep=0pt, scale=  0.60] at ( 21.60,150.50) {4000};

\node[text=drawColor,anchor=base east,inner sep=0pt, outer sep=0pt, scale=  0.60] at ( 21.60,179.97) {5000};
\end{scope}
\begin{scope}
\path[clip] (  0.00,  0.00) rectangle (332.44,195.13);
\definecolor[named]{drawColor}{rgb}{1.00,0.00,0.00}

\path[draw=drawColor,line width= 0.4pt,dash pattern=on 4pt off 4pt ,line join=round,line cap=round] ( 38.71, 44.27) --
	(100.67, 52.20) --
	(162.62, 70.35) --
	(224.58, 94.51) --
	(286.53,158.46);

\path[draw=drawColor,line width= 0.4pt,line join=round,line cap=round] ( 37.52, 43.07) rectangle ( 39.91, 45.47);

\path[draw=drawColor,line width= 0.4pt,line join=round,line cap=round] ( 99.47, 51.00) rectangle (101.86, 53.39);

\path[draw=drawColor,line width= 0.4pt,line join=round,line cap=round] (161.42, 69.15) rectangle (163.82, 71.55);

\path[draw=drawColor,line width= 0.4pt,line join=round,line cap=round] (223.38, 93.32) rectangle (225.77, 95.71);

\path[draw=drawColor,line width= 0.4pt,line join=round,line cap=round] (285.33,157.26) rectangle (287.73,159.66);
\definecolor[named]{drawColor}{rgb}{0.00,0.00,0.00}

\node[text=drawColor,anchor=base,inner sep=0pt, outer sep=0pt, scale=  0.60] at (162.62,  1.44) {Files in Directory};

\node[text=drawColor,rotate= 90.00,anchor=base,inner sep=0pt, outer sep=0pt, scale=  0.60] at (  5.76,108.36) {Time (Milliseconds)};

\path[draw=drawColor,line width= 0.4pt,line join=round,line cap=round] ( 38.71,182.04) rectangle ( 94.12,160.44);
\definecolor[named]{drawColor}{rgb}{0.00,0.00,1.00}

\path[draw=drawColor,line width= 0.4pt,line join=round,line cap=round] ( 40.33,174.84) -- ( 51.13,174.84);
\definecolor[named]{drawColor}{rgb}{1.00,0.00,0.00}

\path[draw=drawColor,line width= 0.4pt,dash pattern=on 4pt off 4pt ,line join=round,line cap=round] ( 40.33,167.64) -- ( 51.13,167.64);
\definecolor[named]{drawColor}{rgb}{0.00,0.00,1.00}

\path[draw=drawColor,line width= 0.4pt,line join=round,line cap=round] ( 45.73,174.84) circle (  1.35);
\definecolor[named]{drawColor}{rgb}{1.00,0.00,0.00}

\path[draw=drawColor,line width= 0.4pt,line join=round,line cap=round] ( 44.54,166.44) rectangle ( 46.93,168.83);
\definecolor[named]{drawColor}{rgb}{0.00,0.00,0.00}

\node[text=drawColor,anchor=base west,inner sep=0pt, outer sep=0pt, scale=  0.60] at ( 56.53,172.77) {Structured};

\node[text=drawColor,anchor=base west,inner sep=0pt, outer sep=0pt, scale=  0.60] at ( 56.53,165.57) {Unstructured};
\end{scope}
\end{tikzpicture}

\end{frame}

\begin{frame}[fragile]
\frametitle{Conclusion}
\begin{itemize}
\item Experimental results look promising
\begin{itemize}
\item Medium-sized collections
\end{itemize}
\item Future Work
\begin{itemize}
\item Extensibility
\item Scalability
\begin{itemize}
\item NDLTD Union Catalog (\~{}2 million records)
\end{itemize}
\item Reference implementation
\begin{itemize}
\item Design principles
\end{itemize}
\end{itemize}
\end{itemize}
\end{frame}

\begin{frame}[fragile]
%\frametitle{Thank You}
\begin{center}
\bigskip
\textbf{\fontsize{18}{18}\selectfont Questions?} \\
\vspace{\stretch{0.5}}
\textbf{\fontsize{18}{18}\selectfont Additional Information} \\
\bigskip
\textbf{\fontsize{18}{18}\selectfont \url{http://dl.cs.uct.ac.za}}
\end{center}
\end{frame}

\end{document}
